\chapter{Correttezza e validazione}
\label{cap:validazione}

La traccia chiede il collaudo funzionale contro un'implementazione seriale
minimale. Il progetto la usa come oracolo, ma la strategia di correttezza \`e
articolata su quattro livelli, perch\'e i modi di sbagliare sono di natura
diversa e un solo controllo non li copre.

\section{L'oracolo seriale}

\file{src/serial/serial.c} \`e la trascrizione letterale della
\cref{eq:definizione}: tre cicli annidati nell'ordine $i \to c \to j$, cio\`e lo
schema~B. La scelta \`e deliberata: l'oracolo deve essere \emph{ovviamente}
corretto, non veloce, e usare un ordine dei cicli diverso da quello dei backend
riduce la probabilit\`a che un errore concettuale sull'ordinamento passi
inosservato in entrambi.

Con \code{--check} il rank 0 raccoglie $Y$ distribuita, ricostruisce da s\'e $A$ e
$X$ globali attraverso il generatore (\cref{sec:generazione}) --- senza
riceverle da nessuno --- calcola l'oracolo e confronta. Il costo \`e $O(MNk)$
seriale, che \`e proprio il motivo per cui la validazione non fa mai parte di una
campagna di misura.

\section{Una tolleranza che dipende da $N$}
\label{sec:tolleranza}

La tolleranza del confronto non \`e una costante, e non pu\`o esserlo: ogni
elemento di $Y$ \`e una riduzione su $N$ termini, e l'errore di arrotondamento di
una riduzione dipende da $N$.

Con dati a media nulla --- che \`e esattamente ci\`o che il generatore produce, e
per questo motivo --- gli arrotondamenti hanno segno indipendente e l'errore
cresce come $\varepsilon\sqrt{N}$, non come il pessimistico $\varepsilon N$ del
caso in cui tutti i termini abbiano lo stesso segno. La costante del modello \`e
\textbf{misurata}, non stimata: confrontando il prodotto distribuito con l'oracolo
su $N = 257 \dots \num{40000}$, in entrambe le precisioni, l'errore relativo in
norma $L_2$ vale
\begin{equation}
  \mathrm{err} \;\approx\; 0{,}17 \cdot \varepsilon \sqrt{N},
\end{equation}
con un rapporto che resta fra $0{,}11$ e $0{,}17$ su tutto l'intervallo. La
soglia adottata \`e quindi
\begin{equation}
  \mathrm{tol}(N) = 8 \cdot \varepsilon \cdot \sqrt{N},
  \label{eq:tolleranza}
\end{equation}
cio\`e un margine uniforme di circa $47\times$ sul modello, su tutte le taglie e
in entrambe le precisioni.

Il confronto con una soglia fissa mostra perch\'e la scelta conta, e sbaglia in
due modi opposti proprio nell'intervallo di taglie che il progetto misura:

\begin{table}[H]
\centering
\footnotesize
\begin{tabular}{@{}llp{0.55\linewidth}@{}}
\toprule
\textbf{Soglia fissa} & \textbf{Difetto} & \textbf{Conseguenza} \\
\midrule
\code{double}, $10^{-12}$ & troppo larga
  & a $N = 257$ l'errore reale \`e $5{,}6\cdot10^{-16}$: la soglia \`e $\sim 1800$
    volte l'errore, e a $N = \num{40000}$ ancora $\sim 140$ volte. Un bug che
    gonfiasse l'errore di due ordini di grandezza passerebbe per PASS. \\
\addlinespace
\code{float}, $10^{-5}$ & si restringe con $N$
  & il margine passa da $33\times$ a $N = 257$ a $2{,}5\times$ a
    $N = \num{40000}$, che \`e la taglia a cui la campagna lavora. Basta un
    backend con un altro ordine di riduzione per produrre un FAIL su un
    risultato corretto --- e siccome il binario esce con stato non nullo, quel
    FAIL abortisce la campagna. \\
\bottomrule
\end{tabular}
\caption{Perch\'e una soglia costante non funziona in nessuna delle due
precisioni.}
\end{table}

Un ultimo dettaglio che conta: $N$ nella \cref{eq:tolleranza} \`e il numero
\emph{globale} di colonne di $A$, non $\nloc$, perch\'e \`e su tutte e $N$ che
l'oracolo somma.

\section{Il test delle funzioni indice}
\label{sec:validazione-indici}

\begin{daverificare}
La cartella \file{test/} non \`e presente nella copia di lavoro corrente, mentre
il \code{Makefile} continua a costruire \file{bin/test\_index} da
\file{test/test\_index.c} e la sonda \code{check-cxx} da
\file{test/cxx\_iface.cpp}: allo stato attuale \code{make}, \code{make test} e
\code{make check-cxx} non possono funzionare. Ripristinare la cartella dal
repository (\code{git checkout -{}- test/}) prima della consegna, oppure
aggiornare questa sezione se i test sono stati rimossi di proposito.
\end{daverificare}

Le tre funzioni di \cref{sec:indici} hanno un test isolato
(\file{test/test\_index.c}, \code{make test}) che non richiede MPI. Verifica su
molte combinazioni $(n, p)$, inclusi i casi limite $p > n$ e $n \bmod p \neq 0$,
tre propriet\`a:
\begin{itemize}[nosep,left=0pt]
  \item $\sum_i \mathrm{size}(n,p,i) = n$: la partizione copre tutto;
  \item $\mathrm{start}(n,p,i+1) = \mathrm{start}(n,p,i) + \mathrm{size}(n,p,i)$:
        le parti sono contigue e non si sovrappongono;
  \item $\mathrm{owner}$ \`e l'inversa di $\mathrm{start}$/$\mathrm{size}$: le due
        direzioni sono coerenti.
\end{itemize}
Sono i controlli che intercettano l'errore di un solo elemento, cio\`e proprio
quello che si propagherebbe in silenzio in tutto il resto del sistema.

\section{La matrice di collaudo}

\begin{table}[H]
\centering
\footnotesize
\caption{I controlli eseguiti e che cosa ciascuno intercetta.}
\label{tab:collaudo}
\begin{tabular}{@{}p{0.34\linewidth}p{0.58\linewidth}@{}}
\toprule
\textbf{Comando} & \textbf{Che cosa protegge} \\
\midrule
\code{make test} & partizionamento degli indici, senza MPI \\
\code{make check-cxx} & \file{kernel.h} e \file{util.h} restano includibili da
  C++ \emph{e} tutti i simboli dell'interfaccia conservano il linkage C
  (verificato con \code{nm}): intercetta le regressioni di \code{extern "C"} e
  \code{restrict} su qualunque macchina, anche senza \code{nvcc} \\
\code{make check} & prodotto distribuito contro oracolo seriale, doppia
  precisione \\
\code{make PREC=float check} & stessa validazione in singola precisione: \`e il
  caso in cui una tolleranza sbagliata si manifesta \\
\code{make check-padding} & $\code{lda} = \nloc + 8$: verifica che tutti i
  backend onorino la leading dimension di $A$ \\
\code{make KERNEL=<nome> check} & la stessa validazione su ciascun backend GPU,
  che ha un ordine di riduzione diverso da quello di CPU \\
\code{make -B KERNEL=cuda\_warp EXTRA\_NVCCFLAGS='-Xptxas -v'} & ispezione
  dell'uso dei registri per warp: non \`e correttezza ma \`e il controllo che dice
  se la specializzazione su $k$ sta facendo ci\`o che si crede \\
\bottomrule
\end{tabular}
\end{table}

La validazione va eseguita su ogni combinazione di backend e precisione
\emph{prima} della campagna corrispondente, su taglie piccole. Un backend
validato solo in doppia precisione non \`e validato: l'ordine della riduzione
cambia con il kernel, e in singola precisione l'effetto \`e misurabile.

\begin{daverificare}
Riportare l'esito effettivo dei controlli sul server, con il valore di
\code{rel\_err} misurato e la soglia corrispondente, per almeno: backend di CPU
in \code{double} e \code{float}; \kernelname{cuda\_warp} e
\kernelname{cuda\_warp\_smem} in entrambe le precisioni; \kernelname{cublas};
il caso con padding. Una tabella con colonne
(backend, precisione, $N$, \code{rel\_err}, \code{tol}, esito) \`e la forma pi\`u
compatta.
\end{daverificare}

\begin{table}[H]
\centering
\footnotesize
\caption{Esiti della validazione. \TODO{compilare dopo l'esecuzione sul server}}
\label{tab:esiti-validazione}
\begin{tabular}{@{}llrrrl@{}}
\toprule
\textbf{Backend} & \textbf{Precisione} & $N$ & \code{rel\_err} & \code{tol} & \textbf{Esito} \\
\midrule
\kernelname{scheme\_a} & double & \dato & \dato & \dato & \dato \\
\kernelname{scheme\_a} & float & \dato & \dato & \dato & \dato \\
\kernelname{scheme\_a} (padding) & double & \dato & \dato & \dato & \dato \\
\kernelname{cuda\_naive} & double & \dato & \dato & \dato & \dato \\
\kernelname{cuda\_warp} & double & \dato & \dato & \dato & \dato \\
\kernelname{cuda\_warp} & float & \dato & \dato & \dato & \dato \\
\kernelname{cuda\_warp\_smem} & double & \dato & \dato & \dato & \dato \\
\kernelname{cuda\_warp\_smem} & float & \dato & \dato & \dato & \dato \\
\kernelname{cuda\_warp\_tiled} & double & \dato & \dato & \dato & \dato \\
\kernelname{cublas} & double & \dato & \dato & \dato & \dato \\
\bottomrule
\end{tabular}
\end{table}
